\section{Diagnóstico do estágio atual}
\label{sec:diagnostico}

\subsection{Fase do negócio}

A \projeto{} encontra-se na fase de \textbf{validação}, em estágio inicial e anterior à tração. O \textit{Startups Report Brasil 2024}, do Observatório Sebrae Startups, mostra que 33,45\% das startups mapeadas estão em validação e que apenas 3,72\% se encontram na etapa de escala \cite{sebrae2025estagios}; a \projeto{} integra, portanto, o grupo mais numeroso do ecossistema, e o aprendizado de validação é seu principal ativo neste momento. A Tabela~\ref{tab:evidencias} reúne as evidências que sustentam esse diagnóstico.

\begin{table}[htb]
\centering
\caption{Evidências do estágio de validação}
\label{tab:evidencias}
\fontsizetable
\begin{tabularx}{\textwidth}{L{3.3cm} Y}
\toprule
\textbf{Dimensão} & \textbf{Situação atual}\\
\midrule
Produto & Protótipo funcional (versão \textit{Alpha 3i}) com cadastro e histórico do animal, calendário de vacinas, prontuário com código QR, triagem orientativa, rede solidária e agendamento; dados locais e demonstrativos, sem pagamentos reais, notificações ou geolocalização \cite{cassino2026repo}. O mapa da rede é esquemático e o login por conta Google é apenas simulado \cite{grupo2026a6}. Há ainda aplicação web temporária e arquivos de instalação para teste.\\
Evidência com usuários & Cinco sessões qualitativas de demonstração guiada com perfis distintos, com recepção favorável e sugestões de melhoria, sobretudo quanto à clareza dos créditos e do código QR \cite{grupo2026a6}. A amostra é pequena e exploratória: não houve execução autônoma de tarefas, e não se afirmam aquisição, retenção, tração comercial ou disposição a pagar das classes C, D e E \cite{grupo2026a9}.\\
Receita e clientes & Sem receita, sem usuários pagantes e sem clínicas contratadas. Os valores de receita são projeções do grupo. O faturamento do setor pet (R\$~77,96 bilhões em 2025 \cite{abempet2026}) contextualiza o segmento, mas não estima a demanda alcançável pela \projeto{}.\\
Formalização & Proposta de constituição como sociedade empresária limitada com quatro sócios \cite{grupo2026a10}; empresa ainda não constituída.\\
Equipe & Quatro integrantes com frentes e substitutos definidos \cite{grupo2026a11}.\\
\bottomrule
\end{tabularx}
\fonteautores
\end{table}

\subsection{Principais desafios para o crescimento}

\begin{table}[htb]
\centering
\caption{Desafios para o crescimento, evidências e implicações}
\label{tab:desafios}
\fontsizetable
\begin{tabularx}{\textwidth}{L{0.7cm} L{3.4cm} Y Y}
\toprule
\textbf{N.º} & \textbf{Desafio} & \textbf{Evidência} & \textbf{Implicação para o plano}\\
\midrule
1 & Disposição a pagar não medida & O teste de aceitação do preço está planejado, mas não realizado \cite{grupo2026a9}. & Validar preço do PetCuida+ e do Plano Parceiro antes de ampliar o investimento em aquisição.\\
2 & Dependência de convênios & Sem os R\$~1.500,00 mensais de convênio e patrocínio, o resultado do cenário enxuto é de R\$~$-$653,68 no mês 12 \cite{grupo2026a9}. & Contratar ao menos um convênio ou patrocínio é condição para a sustentabilidade inicial.\\
3 & Fundo de liquidação de créditos & O fundo custa R\$~6,60 por atendimento, valor superior à margem de R\$~5,24 do atendimento antes do fundo; a margem após o fundo é de R\$~$-$1,36 \cite{grupo2026a9}. & Reduzir o fundo (patrocínio de campanhas, limites de uso) é pré-requisito da escala.\\
4 & Custo de aquisição não medido & O CAC não foi observado; os R\$~150,00 mensais de divulgação não comprovam CAC baixo \cite{grupo2026a9}. & Medir o CAC no piloto e adotar um teto de gasto por usuário (Subseção~\ref{sub:cac}).\\
5 & Arranque de dois lados & Tutores só valorizam a plataforma se houver clínicas, e clínicas só pagam se houver tutores \cite{rochet2003,eisenmann2006}. & Concentrar a oferta numa praça pequena e subsidiar o lado mais sensível a preço.\\
6 & Segurança e dados pessoais & O protótipo ainda não está pronto para dados reais; a LGPD exige base legal e consentimento \cite{brasil2018lgpd}. & Endurecer a segurança e formalizar privacidade antes do piloto com usuários reais.\\
7 & Oferta local não comprovada & Não há clínicas ou organizações contratadas; horários, preços sociais e critérios de encaminhamento ainda não foram validados com parceiros reais. & Descobrir e validar a oferta local no piloto, antes de prometer encaminhamentos ao tutor.\\
8 & Capacidade da equipe & Quatro integrantes com outras obrigações; conhecimento técnico concentrado \cite{grupo2026a11}. & Limitar o número de frentes simultâneas e manter revisão em dupla.\\
\bottomrule
\end{tabularx}
\fonteautores
\end{table}
